% Part: many-valued-logic
% Chapter: three-valued-logics
% Section: multiple-designation
%READERNOTE{SOL6-A142: फक्त निर्दिष्ट मूल्ये बदलून लूकासेविच मॅट्रिक्सला R-Mingle म्हणणे अयोग्य होते. R-Mingle चे वेगळे सशर्त कोष्टक Kooi आणि Tamminga यांच्या primary paper मधील विभाग 4, पृ. 10 शी जुळवले; जोडलेल्या असत्यता-स्थिराचे मूल्य F स्पष्ट केले. LP आणि Hallden यांच्या classical tautologies ची समानता केवळ समान चार-संयोजकीय खंडात आहे.}

\documentclass[../../../include/open-logic-section]{subfiles}

\begin{document}

\olfileid{mvl}{thr}{mul}

\olsection{$\True$ व्यतिरिक्तही मूल्य निर्दिष्ट करणे}

आतापर्यंत पाहिलेल्या सर्व तर्कशास्त्रांत निर्दिष्ट
सत्यतामूल्यांचा संच $V^+ = \{\True\}$ होता; म्हणजे एखाद्या
विधान निर्दिष्ट मानले जाई, जर आणि फक्त जर त्याचे
सत्यतामूल्य~$\True$ असेल.
पण एखादे विधान फक्त~$\False$ नसले, तरी ते सत्य मानता येईल.
अशा वेळी, उदाहरणार्थ, मॅट्रिक्समध्ये
$V^+ = \{\True, \Undef\}$ असे ठरवून तर्कशास्त्र मिळेल.

\begin{defn}
\emph{विरोधापत्तीचे तर्कशास्त्र}~$\LogLP$ पुढील मॅट्रिक्सने
परिभाषित होते:
\begin{enumerate}
  \item मानक विधानीय भाषा $\Lang L_0$ चा केवळ
  $\lnot$, $\land$, $\lor$, $\lif$ हे संयोजक असलेला खंड.
  \item सत्यतामूल्यांचा संच $V = \{\True, \Undef, \False\}$.
  \item $\True$ आणि $\Undef$ निर्दिष्ट आहेत, म्हणजे $V^+ = \{\True, \Undef\}$.
  \item सत्यता-फलने प्रबळ क्लिनी तर्कशास्त्रातील फलनांसारखीच आहेत.
\end{enumerate}
\end{defn}

\begin{defn}
Halld\'en यांचे \emph{निरर्थकतेचे तर्कशास्त्र}~$\LogHal$
पुढील मॅट्रिक्सने परिभाषित होते:
\begin{enumerate}
  \item मानक विधानीय भाषा $\Lang L_0$ च्या
  $\lnot$, $\land$, $\lor$, $\lif$ या संयोजकांच्या खंडाचा
  $1$-स्थानी संयोजक~$+$ जोडून केलेला विस्तार.
  \item सत्यतामूल्यांचा संच $V = \{\True, \Undef, \False\}$.
  \item $\True$ आणि $\Undef$ निर्दिष्ट आहेत, म्हणजे $V^+ = \{\True, \Undef\}$.
  \item सत्यता-फलने दुर्बल क्लिनी तर्कशास्त्रातील फलनांसारखीच
  आहेत; त्याखेरीज ``अर्थहीन आहे'' हा परिकर्मी आहे:
  \begin{center}
    \begin{tabular}{c|c}
    $\tf{+}$ & \\
    \hline
    $\True$ & $\False$ \\
    $\Undef$ & $\True$ \\
    $\False$ & $\False$
    \end{tabular}
  \end{center}
\end{enumerate}
\end{defn}

यांची सत्यताकोष्टके क्लिनी तर्कशास्त्रांशी जुळतात; पण
त्यांच्या विपरीत यांत सर्वतः सत्य सूत्रे \emph{आहेत}.

\begin{prop}\ollabel{prop:LP-taut-CL} $\LogLP$ मधील सर्वतः सत्य
  सूत्रे आणि अभिजात विधानीय तर्कशास्त्रातील सर्वतः सत्य सूत्रे
  एकच आहेत.
\end{prop}

\begin{proof}
  \olref[syn][sub]{prop:mvl-cl} नुसार $\Entails[\LogLP] !A$
  असेल, तर $\Entails[\LogCL] !A$. उलट दिशा दाखवण्यासाठी
  आपण सिद्ध करू की $\pValue v(!A)[\LogKs] = \False$ असेल असे
  !!a{valuation}~$\pAssign v\colon \PVar \to \{\False, \True,
  \Undef\}$ असल्यास $\pValue {v'}(!A)[\LogCL] = \False$
  असे !!a{valuation}~$\pAssign {v'}\colon \PVar \to \{\False, \True\}$
  असते. यावरून $\LogLP$ साठी अपेक्षित निष्कर्ष मिळतो. कारण
  $\LogKs$ आणि $\LogLP$ यांची वैशिष्ट्यपूर्ण सत्यता-फलने
  सारखीच आहेत आणि $\LogLP$ मध्ये फक्त $\False$ हेच मूल्य
  निर्दिष्ट नाही ($\LogLP$ आणि~$\LogKs$ यांतील तोच एक फरक
  आहे). म्हणून $\Entails/[\LogLP] !A$ असेल, तर कोणत्यातरी
  !!{valuation}~$\pAssign v$ साठी $\pValue v(!A)[\LogLP] = \pValue
  v[\LogKs](!A) = \False$. सिद्ध करावयाच्या विधानानुसार
  $\pValue{v'}[\LogCL](!A) = \False$, म्हणजे $\Entails/[\LogCL] !A$.

  हे सिद्ध करण्यासाठी आधी $\pAssign {v'}$ ची व्याख्या अशी करू:
  \[
    \pAssign {v'}(p) =
    \begin{cases}
    \True & \text{जर } \pAssign {v}(p) \in \{\True, \Undef\}\\
    \False & \text{इतर वेळी}
  \end{cases}
  \]
  आता $!A$ वरील विगमनाने दाखवू की (a)~$\pValue v(!A)[\LogKs]
  = \False$ असेल, तर $\pValue {v'}(!A)[\LogCL] = \False$;
  आणि (b)~$\pValue v(!A)[\LogKs] = \True$ असेल, तर
  $\pValue {v'}(!A)[\LogCL] = \True$.
  \begin{enumerate}
    \item विगमनाचा आधार: $!A \ident p$. (a) किंवा (b) चा
    पूर्वपक्ष लागू असेल, तेव्हा मूळ मूल्य अनुक्रमे असत्य किंवा
    सत्य असते; त्या प्रकरणांतच
    \olref[syn][val]{defn:pValue} नुसार $\pValue v(!A)[\LogKs]
    = \pAssign v(p) = \pValue {v'}(!A)[\LogCL]$ ही समानता
    मिळते, आणि त्यावरून (a) व~(b) दोन्ही सिद्ध होतात.

    विगमनाच्या पायरीसाठी पुढील प्रकरणे पाहू:
    \item $!A \ident \lnot !B$.
    \begin{enumerate}
      \item $\pValue v(\lnot!B)[\LogKs] = \False$ असे समजा.
      $\tf{\lnot}[\LogKs]$ च्या व्याख्येनुसार
      $\pValue v(!B)[\LogKs]  = \True$. विगमन गृहितकातील
      प्रकरण~(b) नुसार $\pValue {v'}(!B)[\LogCL]  = \True$;
      म्हणून $\pValue {v'}(\lnot !B)[\LogCL] = \False$.
      \item $\pValue v(\lnot!B)[\LogKs] = \True$ असे समजा.
      $\tf{\lnot}[\LogKs]$ च्या व्याख्येनुसार
      $\pValue v(!B)[\LogKs]  = \False$. विगमन गृहितकातील
      प्रकरण~(a) नुसार $\pValue {v'}(!B)[\LogCL]  = \False$;
      म्हणून $\pValue {v'}(\lnot !B)[\LogCL] = \True$.
    \end{enumerate}

    \item $!A \ident (!B \land !C)$.
    \begin{enumerate}
      \item $\pValue v(!B \land !C)[\LogKs] = \False$ असे समजा.
      $\tf{\land}[\LogKs]$ च्या व्याख्येनुसार
      $\pValue v(!B)[\LogKs]  = \False$ किंवा
      $\pValue v(!C)[\LogKs]  = \False$. विगमन गृहितकातील
      प्रकरण~(a) नुसार $\pValue {v'}(!B)[\LogCL]  = \False$
      किंवा $\pValue {v'}(!C)[\LogCL]  = \False$;
      म्हणून $\pValue {v'}(!B \land !C)[\LogCL] = \False$.
      \item $\pValue v(!B \land !C)[\LogKs] = \True$ असे समजा.
      $\tf{\land}[\LogKs]$ च्या व्याख्येनुसार
      $\pValue v(!B)[\LogKs]  = \True$ आणि
      $\pValue v(!C)[\LogKs]  = \True$. विगमन गृहितकातील
      प्रकरण~(b) नुसार $\pValue {v'}(!B)[\LogCL]  = \True$
      आणि $\pValue {v'}(!C)[\LogCL]  = \True$;
      म्हणून $\pValue {v'}(!B \land !C)[\LogCL] = \True$.
    \end{enumerate}
  \end{enumerate}

  उरलेली दोन प्रकरणे याच प्रकारची आहेत आणि सरावासाठी
  सोडली आहेत. दुसरा मार्ग असा: वरील पुराव्याने फक्त $\lnot$
  आणि~$\land$ असलेल्या सर्व !!{formula}s साठी निष्कर्ष
  सिद्ध केला आहे. आता $\LogKs$ आणि $\LogCL$ या दोन्हींमध्ये
  कोणत्याही $\pAssign v$ साठी $\pValue v(!B \lor
  !C) = \pValue v(\lnot(\lnot!B \land \lnot !C))$ आणि $\pValue v(!B \lif
  !C) = \pValue v(\lnot(!B \land \lnot !C))$ या समानता
  वापरता येतात.
\end{proof}

\begin{prob}
  \olref[mvl][thr][mul]{prop:LP-taut-CL} चा पुरावा पूर्ण
  करा; म्हणजे $!A \ident (!B \lor !C)$ आणि
  $!A \ident (!B \lif !C)$ या प्रकरणांसाठी (a) व~(b)
  सिद्ध करा.
\end{prob}

\begin{prob}
प्रत्येक अभिजात सर्वतः सत्य सूत्र $\LogHal$ मध्येही
सर्वतः सत्य आहे हे सिद्ध करा.
\end{prob}

दोन्ही तर्कशास्त्रांच्या समान चार-संयोजकीय भाषाखंडात त्यांची
सर्वतः सत्य सूत्रे अभिजात तर्कशास्त्रासारखीच असली, तरी त्यांचे
तार्किक निष्पन्नता संबंध वेगळे आहेत. Halld\'en यांच्या अतिरिक्त
संयोजकाचा समावेश असलेल्या सूत्रांसाठी हे समानतेचे विधान नाही. उदाहरणार्थ,
$\LogLP$ हे \emph{व्याघातसहिष्णु} आहे: त्यात
$\lnot p, p \Entails/ q$. म्हणून
$\lnot !A, !A \Entails !B$ हे स्फोटाचे तत्त्व सर्वसाधारणपणे
लागू होत नाही. ($!A$ आणि $!B$ यांच्या काही प्रकरणांत
ते लागू होते; उदाहरणार्थ, $!B$~हे सर्वतः सत्य सूत्र असेल तर.)

\begin{prob}
  पुढीलपैकी कोणते संबंध (a)~$\LogLP$ आणि (b)~$\LogHal$
  मध्ये खरे ठरतात? प्रत्येकासाठी सत्यताकोष्टक द्या.
  \begin{enumerate}
    \item $p, p \lif q \Entails q$
    \item $\lnot q, p \lif q \Entails \lnot p$
    \item $p \lor q, \lnot p \Entails q$
    \item $\lnot p, p \Entails q$
    \item $p \Entails p \lor q$
    \item $p \lif q, q\lif r \Entails p \lif r$
  \end{enumerate}
\end{prob}

$\LogLuk[3]$ मध्ये $\Undef$ हे मूल्यही निर्दिष्ट मानल्यास
काय होईल?
केवळ निर्दिष्ट मूल्यांचा हा बदल केल्याने पुढील R-Mingle
मॅट्रिक्स मिळत नाही; त्यासाठी सशर्त विधानाचे सत्यता-फलनही बदलते.

\begin{defn}
  \emph{त्रिमूल्य R-Mingle}~$\LogRM[3]$ चा असत्यता-स्थिरासह
  विस्तार पुढील मॅट्रिक्सने परिभाषित होतो:
  \begin{enumerate}
    \item $\lfalse$, $\lnot$, $\land$, $\lor$, $\lif$ असलेली मानक विधानीय भाषा $\Lang L_0$.
    \item सत्यतामूल्यांचा संच $V = \{\True, \Undef, \False\}$.
    \item $\True$ आणि $\Undef$ निर्दिष्ट आहेत, म्हणजे $V^+ = \{\True, \Undef\}$.
    \item $\tf{\lfalse}=\False$. $\lnot$, $\land$ आणि $\lor$
    यांची सत्यता-फलने \L ukasiewicz तर्कशास्त्र~$\LogLuk[3]$
    मधील फलनांसारखीच आहेत. सशर्त विधानाचे सत्यताकोष्टक असे आहे:
    \begin{center}
      \begin{tabular}{c|ccc}
        $\tf{\lif}[\LogRM[3]]$ & $\True$ & $\Undef$ & $\False$ \\
        \hline
        $\True$ & $\True$ & $\False$ & $\False$ \\
        $\Undef$ & $\True$ & $\Undef$ & $\False$ \\
        $\False$ & $\True$ & $\True$ & $\True$
      \end{tabular}
    \end{center}
  \end{enumerate}
  येथे $\Undef$ हे मधले मूल्य सत्य आणि असत्य दोन्ही असल्याचे
  दर्शवते. सशर्त विधानाचे कोष्टक \L ukasiewicz यांच्या कोष्टकापेक्षा
  वेगळे आहे.\footnote{Barteld Kooi आणि Allard Tamminga,
  ``Completeness via Correspondence for Extensions of the Logic of
  Paradox,'' (2012), प्रथम ऑनलाइन आवृत्ती पृ. 10, विभाग 4;
  \url{https://doi.org/10.1017/S1755020312000196}.}
\end{defn}

\begin{prob}
  पुढीलपैकी कोणते संबंध $\LogRM[3]$ मध्ये खरे ठरतात?
  \begin{enumerate}
    \item $p, p \lif q \Entails q$
    \item $p \lor q, \lnot p \Entails q$
    \item $\lnot p, p \Entails q$
    \item $p \Entails p \lor q$
  \end{enumerate}
\end{prob}

निर्दिष्ट मूल्ये बदलल्यामुळे वेगवेगळी सत्यताकोष्टके कधी कधी
एकच तर्कशास्त्र, म्हणजे एकच तार्किक निष्पन्नता संबंध, देऊ
शकतात. उदाहरणार्थ, G\"odel तर्कशास्त्रात~$\{\True\}$ ऐवजी
$V^+ = \{\True, \Undef\}$ घेतल्यास असे घडते.

\begin{prop}\ollabel{prop:gl-udes}
  $V = \{\False, \Undef,\True\}$,
  $V^+=\{\True, \Undef\}$ आणि $3$-मूल्य G\"odel
  तर्कशास्त्राची सत्यता-फलने असलेला मॅट्रिक्स अभिजात
  तर्कशास्त्र परिभाषित करतो.
\end{prop}

\begin{proof}
  सराव.
\end{proof}

\begin{prob}
  G\"odel तर्कशास्त्राप्रमाणेच पण
  $V^+=\{\True,\Undef\}$ घेऊन परिभाषित केलेल्या
  तर्कशास्त्र~$\Log L$ साठी $\Gamma \Entails/[\Log L] !B$
  असेल, तर $\Gamma \Entails/[\LogCL] !B$ हे दाखवून
  \olref[mvl][thr][mul]{prop:gl-udes} सिद्ध करा.
  \olref[mvl][thr][mul]{prop:LP-taut-CL} मधील कल्पना
  वापरा; मात्र गुणधर्म (a) आणि~(b) सिद्ध करण्याऐवजी
  $\pValue v(!A)[\LogGod] = \False$ जर आणि फक्त जर
  $\pValue {v'}(!A)[\LogCL] = \False$ हे दाखवा (आणि म्हणून
  $\pValue v(!A)[\LogGod] \in \{\True,\Undef\}$ जर आणि
  फक्त जर $\pValue {v'}(!A)[\LogCL] = \True$).
  यावरून विधान का सिद्ध होते ते स्पष्ट करा.
\end{prob}

\end{document}
